\begin{figure}[h!]
\centering
\resizebox{0.96\columnwidth}{!}{%
\begin{tikzpicture}[>=Latex, thick,
  lvl/.style={rectangle, rounded corners, draw=black!60, fill=black!3,
              text width=48mm, align=left, inner sep=4pt, font=\footnotesize},
  reg/.style={rectangle, draw=black, fill=white, inner sep=3pt,
              font=\scriptsize, align=center, minimum height=6mm},
  fix/.style={reg, fill=black!8},
  lab/.style={font=\tiny, align=center}]

  % --- the two levels, described ---
  \node[lvl] (l1) {\textbf{Level 1: the total fleet} $X=\sum_{g,t}n_{gt}$\\[2pt]
     transport cost is exactly $cK\cdot X$, so a bound on the total has
     nothing to compensate it --- unlike a per-period bound, which the
     relaxation answers by shipping in the neighbouring period};
  \node[lvl, below=6mm of l1] (l2) {\textbf{Level 2: the fleet vector, inside $X=X^{\star}$}\\[2pt]
     order the pairs $(g,t)$; at the first pair where a plan may differ
     from the incumbent branch $<$, $=$, $>$ and recurse on the equality
     (Murty). Every fleet vector differs from $\mathbf{n}^{\star}$ at a
     first index, so the regions are disjoint and exhaustive};
  \node[lvl, below=6mm of l2] (l3) {\textbf{At every node}\\[2pt]
     a plan is \emph{constructed} inside the region and its trucks join
     the shared pool; the children partition around that plan. The tree
     finds columns; the construction finds plans};

  % --- the tree itself ---
  \begin{scope}[xshift=70mm, yshift=4mm]
    \node[reg] (r) at (0,0) {root: incumbent $\mathbf{n}^{\star}$, $X^{\star}$};
    \node[reg] (a) at (-3.0,-1.5) {$X\ge X^{\star}+1$};
    \node[fix] (b) at (0,-1.5) {$X=X^{\star}$};
    \node[reg] (c) at (3.0,-1.5) {$X\le X^{\star}-1$};
    \draw[->] (r) -- (a); \draw[->] (r) -- (b); \draw[->] (r) -- (c);
    \node[lab, below=1pt of a, text width=26mm] {a whole truck with no offset:\\ usually pruned by bound};
    \node[lab, below=1pt of c, text width=26mm] {the only region that can beat\\ the incumbent outright: certified};

    \node[reg] (b1) at (-2.6,-3.4) {$n_{1}\le n^{\star}_{1}-1$};
    \node[fix] (b2) at (0,-3.4) {$n_{1}=n^{\star}_{1}$};
    \node[reg] (b3) at (2.6,-3.4) {$n_{1}\ge n^{\star}_{1}+1$};
    \draw[->] (b) -- (b1); \draw[->] (b) -- (b2); \draw[->] (b) -- (b3);

    \node[reg] (c1) at (-2.6,-5.0) {$n_{2}\le n^{\star}_{2}-1$};
    \node[fix] (c2) at (0,-5.0) {$n_{2}=n^{\star}_{2}$};
    \node[reg] (c3) at (2.6,-5.0) {$n_{2}\ge n^{\star}_{2}+1$};
    \draw[->] (b2) -- (c1); \draw[->] (b2) -- (c2); \draw[->] (b2) -- (c3);
    \node[lab, below=1pt of c2] {$\vdots$ \ to a fixed depth};
    \node[lab, anchor=west, text width=30mm] at (3.9,-4.2)
      {shaded: transport fixed;\\ everything inside is\\ inventory allocation};
  \end{scope}
\end{tikzpicture}}
\caption{The two levels of the fleet tree. Level 1 partitions on the
total fleet $X$, whose cost $cK\cdot X$ the relaxation cannot offset;
level 2 partitions the region $X=X^{\star}$ on the fleet vector itself,
one pair at a time around the plan constructed at each node. Bounds on
$X$ and on $n_{gt}$ are rows of the master that already carry duals,
so the pricing subproblem is untouched, and a column found in any
region is valid in every other.}
\label{fig:branching}
\end{figure}
